% Appendix B, typeset from info/quartus_workflow.md.

\chapter{Quartus Prime Lite and the DE0-CV}
\label{app:quartus}

This is the book's tool reference for synthesis and hardware. Through \chapref{c1:ch} to
\ref{c18:ch} it is \textbf{teacher-facing}: the Quartus and DE0-CV work is demonstrated from the
front, starting with the brake assistant in \chapref{c1:ch}, and no participant needs Quartus or a
board of their own.

\begin{aside}
\Chapref{c19:ch} is the exception, and it is the one that counts. The bring-up is demonstrated first
and then handed over: every group synthesises and programs its own \code{can_spi_node} and
demonstrates its own two-node link. So read \secref{app:quartus:project} to \ref{app:quartus:program}
as participant-facing \textemdash{} you need \secref{app:quartus:project} and
\ref{app:quartus:sources} to create the project and add the source files, and
\secref{app:quartus:compile} is the compilation producing the \file{.sof} that
\secref{app:quartus:program} programs \textemdash{} and read \secref{app:quartus:linux} \emph{before}
you connect anything, since that is the section deciding whether the Programmer can see the board at
all.
\end{aside}

Nothing that is \textbf{assessed} requires the hardware (appendix~\ref{app:project}). The bring-up is
assessed on the simulation behind it, not on whether the wire worked on that particular day.
Everything described here ends on real FPGA hardware, verified by hand. That is a different layer from
the exercises, which you verify on your own laptop with a self-checking testbench under GHDL
(appendix~\ref{app:sim}).

% ------------------------------------------------------------------------------------------
\section{What you need}
\label{app:quartus:need}

\begin{itemize}
  \item \textbf{Quartus Prime Lite} (the free edition): the manufacturer's FPGA tool, covering
    synthesis, place and route (``fitting''), and programming file generation.
  \item \textbf{Cyclone V device support}: a separate add-on package. The DE0-CV's FPGA,
    \code{5CEBA4F23C7N}, is a Cyclone V device, and Quartus does not install every family's support
    by default.
  \item \textbf{USB-Blaster driver}: the DE0-CV is programmed over USB using the USB-Blaster
    protocol. The driver usually installs alongside Quartus, but sometimes needs a manual nudge
    (\secref{app:quartus:trouble}).
  \item A \textbf{Terasic DE0-CV} board and a USB cable.
\end{itemize}

Quartus Prime Lite exists for both Windows and Linux, and the design steps below are the same on
either host. \textbf{Getting the board visible to the Programmer is not}, and that is where every
host-specific problem in this course lives: \secref{app:quartus:linux} covers Linux and WSL,
\secref{app:quartus:trouble} covers Windows.

\begin{warning}
A recommendation straight away, since it saves an evening: the course's GHDL toolchain runs on WSL or
Ubuntu, but \textbf{WSL2 cannot see a USB device at all without extra work}
(\secref{app:quartus:linux}). If you are on Windows, install Quartus natively on Windows and keep WSL
for GHDL. The two never need to see each other; they only share the \file{.vhd} files, which live in
the Windows file system anyway.
\end{warning}

% ------------------------------------------------------------------------------------------
\section{Installing Quartus Prime Lite}
\label{app:quartus:install}

\begin{enumerate}
  \item Download \textbf{Quartus Prime Lite Edition} from the manufacturer's download portal. Search
    for ``Quartus Prime Lite download'': Intel's FPGA business became \textbf{Altera} again in 2024,
    so the pages have moved and older deep links rot. What is stable to search for is the product
    name, not the URL.
    \begin{itemize}
      \item Choose the \textbf{Lite} edition, which is the free one and needs no licence file.
        Standard and Pro do not cover Cyclone V for free.
      \item Any reasonably recent Lite edition works; the course does not depend on a particular
        point release.
      \item Pick the Windows or Linux installer according to where you intend to run it, and read the
        warning above if you are on Windows with WSL.
    \end{itemize}
  \item During installation, make sure \textbf{Cyclone V} is selected among the device families. If
    Quartus is already installed without it, add the Cyclone V support afterwards, either by running
    the installer again or via Quartus's own ``Additional Software'' installer.
  \item When the DE0-CV board is connected over USB for the first time, Windows should detect the
    USB-Blaster and either install the driver automatically or ask you to point out Quartus's own
    driver files (usually under Quartus's installation directory). If Windows finds no driver
    automatically, see \secref{app:quartus:trouble}.
\end{enumerate}

\subsection{Two pitfalls in Windows 11}
\label{app:quartus:win11}

A couple of Windows 11's security defaults can interfere with the installation, or with programming
the board later:

\begin{itemize}
  \item \textbf{If the installer does not even open:} find \code{Smart App Control} in Windows
    Security (search for it in the start menu), turn it off and restart.
  \item \textbf{If everything installs cleanly but the board still cannot be programmed}
    (\secref{app:quartus:program}): go to \code{Windows Security → Device security → Core isolation}
    and turn off \textbf{Memory integrity}.
\end{itemize}

\lecfig[0.78]{info/images/mem_integrity.png}{\code{Memory integrity} under Windows 11's core
isolation settings.}{app:quartus:fig:memint}

Neither setting touches Quartus itself; both are general hardening defaults in Windows 11 that happen
to be strict enough to interfere with driver-level USB access.

% ------------------------------------------------------------------------------------------
\section{Creating a new project}
\label{app:quartus:project}

\begin{enumerate}
  \item Open Quartus and start the \textbf{New Project Wizard} (\code{File → New Project Wizard}).
  \item Choose a working directory and a project name. Quartus uses the project name as the default
    name of the top-level entity too, so let it match your VHDL entity's name to save overwriting it
    later (for \file{or_gate.vhd}: name the project \code{or_gate}).
  \item Choose \textbf{Empty Project} as the project type. The course uses neither Quartus's IP
    catalog nor any project template.
  \item Skip adding files at this step (\secref{app:quartus:sources} covers that), unless you already
    know exactly which file it is; either is fine.
  \item On the device selection page, choose exactly the device on the DE0-CV: filter on the
    \textbf{Cyclone V} family, then find and select the device \code{5CEBA4F23C7N}.
  \item Skip the EDA tools page (leave it at ``none''); the course runs its simulations directly in
    GHDL rather than through Quartus.
  \item Finish the wizard.
\end{enumerate}

% ------------------------------------------------------------------------------------------
\section{Adding your VHDL}
\label{app:quartus:sources}

\begin{itemize}
  \item Add the source file: \code{Project → Add/Remove Files in Project...}, browse to your
    \file{.vhd} file and add it.
  \item If the design has more than one file (a top-level entity plus one or more subcomponents), add
    all of them. Quartus resolves the instantiation hierarchy automatically as long as every file is
    in the project.
  \item Confirm that the top-level entity is set correctly: \code{Project → Set as Top-Level Entity},
    with your main file selected. For a single-file design that happens automatically.
\end{itemize}

% ------------------------------------------------------------------------------------------
\section{Assigning pins in the Pin Planner}
\label{app:quartus:pins}

Every port in your entity needs a physical pin on the FPGA before Quartus can produce a working
programming file; otherwise there is no path from the design's ports to a real switch or LED on the
board.

\begin{enumerate}
  \item Run \code{Processing → Start → Start Analysis \& Elaboration} first, then open the Pin
    Planner via \code{Assignments → Pin Planner}. In a brand new project the port list is empty until
    Quartus has elaborated the design and knows its ports.
  \item Each of the entity's ports appears as a row. Fill in the \textbf{Location} column with the pin
    designation of the switch or LED you want to use.
  \item \textbf{A switch or an LED is a free choice. A clock is not.} For \code{SW}, \code{KEY} and
    \code{LEDR}, any convenient ones will do: the design does not care which switch drives which
    input, only that every port reaches one. \code{CLOCK_50} is a specific pad on the device and
    cannot be moved. On the DE0-CV it is \textbf{\code{PIN_M9}}, the 50 MHz oscillator. Every design
    from \chapref{c3:ch} onwards is clocked, so from that chapter this is the one pin number worth
    knowing without looking it up. Confirm it against your own board's files (step 4) before you rely
    on it.
  \item \textbf{Do not copy the pin table out by hand.} Ten \code{LEDR} pins transcribed from a PDF
    are ten chances to transpose two digits, and the failure it gives is the nastiest one in
    \secref{app:quartus:trouble}: a clean compilation that behaves incomprehensibly.
    \begin{itemize}
      \item The DE0-CV's system CD contains a \file{DE0_CV.qsf} where every board pin is already named
        and assigned.
      \item \code{Assignments → Import Assignments…}, point out that file, and every
        \code{CLOCK_50}, \code{SW[n]}, \code{KEY[n]}, \code{LEDR[n]} and \code{GPIO_0[n]} assignment
        lands at once.
      \item Name your entity's ports to match those signal names and there is nothing left to assign.
        That is why \chapref{c19:ch}'s bring-up names its ports \code{CLOCK_50}, \code{KEY},
        \code{SW}, \code{LEDR} and \code{GPIO_0} rather than something prettier.
      \item The Pin Planner (and the board's pin tables in the user manual) is still the way to check
        an individual pin, or assign one the \file{.qsf} file does not cover.
    \end{itemize}
  \item Set the \textbf{I/O Standard} column if Quartus does not infer a sensible default for the pin
    you chose. The DE0-CV's general I/O is 3.3 V LVTTL/LVCMOS. Importing the board's \file{.qsf} (step
    4) sets this too, which is one more reason to prefer that route.
  \item Save the assignments (they are written into the project's \file{.qsf} file). Every change here
    requires recompilation (\secref{app:quartus:compile}) before it takes effect; Quartus does not
    retrofit a \file{.sof} compiled before the change.
\end{enumerate}

\subsection{A clock needs a timing constraint too}
\label{app:quartus:sdc}

Assigning \code{CLOCK_50} a pin tells Quartus \emph{where} the clock comes in. It does not say how
fast it runs, and until you say so the timing analysis has nothing to check against and reports an
unconstrained design. Add a \file{.sdc} file to the project with one line:

\begin{console}
create_clock -name CLOCK_50 -period 20.000 [get_ports CLOCK_50]
\end{console}

20 ns is 50 MHz. With that in place, the timing analysis in \code{Processing → Start Compilation}
becomes meaningful and the report gives you a real Fmax, which is the number \secref{c5:sec:fmax}
sends you here to read. Without it, Fmax is missing, or not to be trusted.

% ------------------------------------------------------------------------------------------
\section{Compiling the design}
\label{app:quartus:compile}

\code{Processing → Start Compilation} (or ▶ in the toolbar) runs the whole flow: analysis and
synthesis, place and route, timing analysis, and assembly into a programming file.

Read the compilation report when it is done, in particular the \textbf{Warnings}: an unconnected port,
an incomplete sensitivity list, or a latch inferred where you meant combinational logic all turn up
there as warnings rather than errors. They are easy to miss if you only check that the tick went
green.

A successful compilation produces a \file{.sof} file (SRAM Object File) in the project's
\file{output_files} directory. That is the one \secref{app:quartus:program} downloads to the board. It
is volatile: it configures the FPGA's SRAM directly and disappears when the power is removed. The
course has no step that burns the configuration permanently.

% ------------------------------------------------------------------------------------------
\section{Programming the DE0-CV}
\label{app:quartus:program}

\begin{enumerate}
  \item Connect the DE0-CV to the computer over USB and switch it on.
  \item Open the Programmer: \code{Tools → Programmer}.
  \item Select the USB-Blaster under \textbf{Hardware Setup}. It should be found automatically if the
    driver installed correctly (otherwise \secref{app:quartus:trouble}).
  \item Make sure the \file{.sof} file from \secref{app:quartus:compile} is in the list and that its
    \textbf{Program/Configure} checkbox is ticked. If it is not there, use \code{Auto Detect} or add
    the file by hand.
  \item Click \textbf{Start}. The progress bar should reach 100 \% and report success.
  \item Test the design directly on the board: toggle the switches you assigned in
    \secref{app:quartus:pins} and confirm that the LED follows the entity's truth table.
\end{enumerate}

% ------------------------------------------------------------------------------------------
\section{Troubleshooting}
\label{app:quartus:trouble}

\begin{description}[style=unboxed,leftmargin=0pt,itemsep=0.8ex]
  \item[The USB-Blaster does not appear in the Programmer.] Check Device Manager in Windows for an
    unknown device or one with a driver error, and manually point out the USB-Blaster driver under
    Quartus's installation directory.
  \item[The board cannot be programmed, with no clear error.] Turn off \code{Memory integrity} under
    \code{Core isolation} on Windows 11 (\secref{app:quartus:win11}). It is by far the most common
    cause once the driver is installed.
  \item[The installer does not start.] Check \code{Smart App Control} under Windows Security
    (\secref{app:quartus:win11}).
  \item[The compilation passes but the design misbehaves on the board.] Check the assignments in the
    Pin Planner (\secref{app:quartus:pins}) \emph{before} you suspect your VHDL. A swapped or missing
    pin assignment gives a board that compiles cleanly and behaves incomprehensibly, since it is not
    wired the way you think.
  \item[The fitter reports unassigned pins as errors.] Some project settings treat that as fatal
    rather than as a warning. Either assign every port in the Pin Planner, or, for a forgotten unused
    port early in development, tell Quartus how to treat the rest:
    \code{Assignments → Device → Device and Pin Options… → Unused Pins}. For the DE0-CV,
    \textbf{``As input tri-stated (with weak pull-up)''} is the safe value. The default in some flows
    drives unused pins low, which on a board where pins reach headers and peripherals is a way of
    fighting another driver. Do not leave ports silently unassigned and expect Quartus to guess.
\end{description}

% ------------------------------------------------------------------------------------------
\section{Programming from Linux and from WSL}
\label{app:quartus:linux}

\Secref{app:quartus:trouble} is Windows. On Linux the driver is no driver at all: the Programmer talks
to the USB-Blaster through libusb, and needs permission to do so.

\subsection*{Linux, natively}
Out of the box the Programmer shows no hardware, since the device node is owned by root. Give it a
udev rule with the lines below in \file{/etc/udev/rules.d/51-usbblaster.rules}:

\begin{console}
# Altera USB-Blaster
SUBSYSTEM=="usb", ATTR{idVendor}=="09fb", ATTR{idProduct}=="6001", MODE="0666"
SUBSYSTEM=="usb", ATTR{idVendor}=="09fb", ATTR{idProduct}=="6002", MODE="0666"
SUBSYSTEM=="usb", ATTR{idVendor}=="09fb", ATTR{idProduct}=="6003", MODE="0666"
# USB-Blaster II
SUBSYSTEM=="usb", ATTR{idVendor}=="09fb", ATTR{idProduct}=="6010", MODE="0666"
SUBSYSTEM=="usb", ATTR{idVendor}=="09fb", ATTR{idProduct}=="6810", MODE="0666"
\end{console}

Reload the rules and then unplug and replug the board:

\begin{shell}
sudo udevadm control --reload-rules && sudo udevadm trigger
\end{shell}

Quartus caches its hardware list in a background daemon, so if the Programmer still shows nothing,
restart it:

\begin{shell}
killall jtagd 2> /dev/null
jtagconfig            # should now list the USB-Blaster and the device chain
\end{shell}

\code{jtagconfig} is the quickest way to tell a permissions problem from a cable problem: if it lists
the board, the hardware is fine and the problem is inside Quartus; if it does not, stop and solve that
before you touch the Programmer.

\subsection*{WSL2}
This is the case the course's own toolchain runs straight into, so it is worth saying outright:
\textbf{WSL2 has no USB passthrough.} The board is invisible to a WSL2 installation however many udev
rules you write, since the kernel never sees the device. There is no setting to change. Two ways out:

\begin{itemize}
  \item \textbf{Run Quartus natively on Windows} (recommended). Keep WSL for GHDL, which needs no
    hardware at all. Your \file{.vhd} files live in the Windows file system and both tools read them.
    That is the route the course assumes, and the one with the fewest moving parts.
  \item \textbf{Forward the device with \code{usbipd-win}}, if you specifically want Quartus inside
    WSL. From an administrator PowerShell on the Windows side:

\begin{console}
winget install usbipd
usbipd list                      # find the USB-Blaster's BUSID
usbipd bind   --busid <BUSID>    # once per device, persists
usbipd attach --wsl --busid <BUSID>
\end{console}

    The device then appears inside WSL and the udev rule above applies. It has to be reattached after
    every unplug and every WSL restart, and Quartus's graphical interface also needs WSLg or an X
    server. Doable, but three more things that can go wrong on a presentation day.
\end{itemize}
